% ============================================================================
% Expérience 01 — Choix des paramètres de clustering
% Rapport complet : stabilité, nselectboot, comparaison
% ============================================================================
% Ton étudiant, simple et clair. Compilation : make (depuis ce dossier)
% ============================================================================

\documentclass[11pt, a4paper]{article}
\usepackage[utf8]{inputenc}
\usepackage[T1]{fontenc}
\usepackage[french]{babel}
\usepackage{amsmath, amssymb}
\usepackage[margin=2.5cm]{geometry}
\usepackage{booktabs}
\usepackage{float}
\usepackage{graphicx}
\usepackage{xcolor}
\usepackage{hyperref}
\hypersetup{colorlinks=true, linkcolor=blue}

\title{Expérience 01 — Comment choisir le nombre de clusters et les paramètres ?}
\author{Pierre Chambet}
\date{Mars 2026}

\begin{document}
\maketitle

\begin{abstract}
On a trois paramètres à choisir : $k$ (nombre de groupes), $\alpha$ et $\omega$
(qui règlent la distance). On compare deux méthodes : la stabilité (notre
protocole) et nselectboot (méthode de la littérature). La stabilité choisit
toujours $k=2$, ce qui pose problème. nselectboot donne des résultats plus
variés et récupère $k=4$ sur Canadian et AEMET, et $k=3$ sur Tecator.
\end{abstract}

\tableofcontents
\newpage

% ============================================================================
\section{Le problème en une phrase}
% ============================================================================

Pour faire du clustering, il faut choisir : combien de groupes ($k$) ? Et
quels paramètres pour la distance ($\alpha$, $\omega$) ? On n'a pas de labels
pour vérifier. Il faut donc des critères qui marchent sans vérité terrain.

% ============================================================================
\section{Les trois paramètres à choisir}
% ============================================================================

\begin{itemize}
  \item \textbf{$k$} : nombre de clusters. Exemple : 2 groupes, 3 groupes, etc.
  \item \textbf{$\alpha$} : mélange entre distance sur les courbes (niveau) et
    sur les dérivées (forme). 0 = niveau seul, 1 = forme seule.
  \item \textbf{$\omega$} : mélange entre partie fonctionnelle (courbes) et
    partie vectorielle (latitude, altitude, etc.). 0 = vectoriel seul, 1 =
    fonctionnel seul.
\end{itemize}

On teste une grille : $k$ de 2 à 6, $\alpha$ et $\omega$ de 0 à 1. On a 4
jeux de données avec des labels (pour évaluer) : Canadian (35 stations, 4
régions), AEMET (73 stations, 4 zones), Growth (93 enfants, 2 sexes), Tecator
(215 échantillons, 3 niveaux de gras).

% ============================================================================
\section{Méthode 1 : La stabilité (notre protocole)}
% ============================================================================

\subsection{L'idée}

Un bon choix de paramètres donne des groupes \textbf{reproductibles} : si on
enlève quelques individus au hasard et qu'on reclustère, on doit retrouver
à peu près les mêmes groupes.

\subsection{Comment on fait}

Pour chaque triplet $(k, \alpha, \omega)$ :
\begin{enumerate}
  \item On fait PAM sur tous les $n$ individus $\to$ partition $\mathcal{C}$.
  \item On répète 150 fois : on garde 80\,\% des individus au hasard, on
    reclustère dessus $\to$ partition $\mathcal{C}_b^*$. On compare
    $\mathcal{C}$ et $\mathcal{C}_b^*$ avec l'ARI (accord entre les deux).
  \item Stabilité = moyenne des 150 ARI.
\end{enumerate}
On garde le triplet avec la stabilité la plus haute.

\subsection{PAM et ARI en bref}

\textbf{PAM} : comme k-means, mais les centres sont de vrais individus
(médoïdes). Ça marche avec n'importe quelle distance.

\textbf{ARI} : mesure l'accord entre deux partitions. Entre $-1$ et $1$. ARI =
1 = tout pareil. On compare toujours sur les mêmes individus (les 80\,\% gardés).

\subsection{Résultat}

La stabilité choisit \textbf{toujours} $k=2$, quel que soit le dataset. Sur
Canadian et AEMET, il y a pourtant 4 classes. On perd de l'information.

% ============================================================================
\section{Méthode 2 : nselectboot (Fang \& Wang, 2012)}
% ============================================================================

\subsection{L'idée}

Cette méthode \textbf{minimise l'instabilité} au lieu de maximiser la
stabilité. L'instabilité mesure : à quel point les clusters changent quand on
refait un tirage bootstrap. Moins ça change, mieux c'est.

\subsection{Comment on fait}

On teste plusieurs couples $(\alpha, \omega)$. Pour chaque couple, on tire des
paires d'échantillons bootstrap, on clustère chacun, on mesure le désaccord.
Le $k$ qui minimise ce désaccord est retenu. Implémentation : fonction
\texttt{nselectboot} du package R \texttt{fpc}.

\subsection{Résultat}

Le $k$ optimal \textbf{change} selon $(\alpha, \omega)$. Sur une grille 21×21
(441 couples), on obtient des $k$ différents selon les zones. En prenant la
médiane sur la grille : Canadian $\to$ 4, AEMET $\to$ 4, Growth $\to$ 2,
Tecator $\to$ 3. Pour Canadian, AEMET et Tecator, c'est le bon $k$.

% ============================================================================
\section{Comparaison des deux méthodes}
% ============================================================================

\begin{table}[H]
\centering
\renewcommand{\arraystretch}{1.3}
\small
\begin{tabular}{lcccccc}
\toprule
\textbf{Dataset} & $k_{\mathrm{vrai}}$ & \multicolumn{2}{c}{\textbf{Stabilité}} & \multicolumn{2}{c}{\textbf{nselectboot}} & \\
& & $k$, $(\alpha,\omega)$ & ARI & $k$, $(\alpha,\omega)$ & ARI & \\
\midrule
Canadian & 4 & 2, (0.30, 0) & 0.54 & 4, (0.25, 0.15) & 0.62 \\
AEMET    & 4 & 2, (0, 1) & 0.27 & 4, (0.35, 0.20) & 0.39 \\
Growth   & 2 & 2, (1, 0.45) & 0.72 & 2, (0.10, 0.05) & 0.42 \\
Tecator  & 3 & 2, (0.40, 0.30) & 0.47 & 3, (0.75, 0.10) & 0.58 \\
\bottomrule
\end{tabular}
\caption{Comparaison : chaque méthode choisit un triplet différent. L'ARI est
calculé sur la partition obtenue avec ce triplet.}
\label{tab:comparaison}
\end{table}

\subsection{Ce qu'on observe}

\begin{itemize}
  \item \textbf{Growth} : les deux méthodes trouvent $k=2$. C'est correct.
    ARI = 0.72, bon résultat.
  \item \textbf{Canadian et AEMET} : la stabilité dit $k=2$, nselectboot dit
    $k=4$ (médiane sur la grille). Le vrai $k$ est 4. nselectboot fait mieux ici.
  \item \textbf{Tecator} : la stabilité dit $k=2$, nselectboot trouve $k=3$
    (le vrai $k$). Avec la grille fine, nselectboot récupère les 3 classes.
\end{itemize}

% ============================================================================
\section{Pourquoi la stabilité favorise $k=2$ ?}
% ============================================================================

Avec $k=2$, on a deux groupes. Deux partitions à deux groupes ont peu de
façons de diverger : soit on est d'accord, soit on sépare différemment. Avec
$k=4$, quatre groupes peuvent se réarranger de plein de façons. Donc l'ARI
entre la partition complète et celle du sous-échantillon est plus bas en
moyenne. La stabilité favorise mécaniquement les petits $k$.

% ============================================================================
\section{Heatmaps et matrices de confusion : deux protocoles, deux résultats}
% ============================================================================

\textbf{Important} : les heatmaps et les matrices de confusion \textbf{dépendent
du protocole}. Ce n'est pas la même chose.

\begin{itemize}
  \item \textbf{Heatmaps stabilité} : montrent la stabilité (moyenne ARI) pour
    chaque $(\alpha, \omega)$ et chaque $k$. Données : \texttt{stabilite\_*.csv}.
  \item \textbf{Heatmaps nselectboot} : montrent l'instabilité pour chaque
    $(\alpha, \omega)$ et chaque $k$. Bleu = instabilité faible (bon), rouge =
    forte (mauvais). Données : \texttt{nselectboot\_*.csv}.
  \item \textbf{Matrices de confusion} : une par méthode, car chaque méthode
    choisit un triplet $(k, \alpha, \omega)$ différent $\to$ partition différente.
\end{itemize}

% ----------------------------------------------------------------------------
\subsection{Matrices de confusion : comparaison côte à côte}
% ----------------------------------------------------------------------------

Lignes = vraies classes. Colonnes = clusters trouvés.

\paragraph{Canadian.} Stabilité : $k=2$, 4 régions en 2 clusters. nselectboot :
$k=4$, les 4 régions sont séparées. ARI passe de 0.54 à 0.62.

\begin{center}
\small
\begin{tabular}{@{}p{4.8cm}p{5.8cm}@{}}
\textbf{Stabilité} ($k=2$, ARI 0.54) &
\textbf{nselectboot} ($k=4$, ARI 0.62) \\[0.5em]
\begin{tabular}{lrr}
\toprule
& C1 & C2 \\
\midrule
Arctic & 0 & 3 \\
Atlantic & 15 & 0 \\
Continental & 1 & 11 \\
Pacific & 1 & 4 \\
\bottomrule
\end{tabular}
&
\begin{tabular}{lrrrr}
\toprule
& C1 & C2 & C3 & C4 \\
\midrule
Arctic & 0 & 0 & 3 & 0 \\
Atlantic & 15 & 0 & 0 & 0 \\
Continental & 1 & 6 & 5 & 0 \\
Pacific & 0 & 4 & 0 & 1 \\
\bottomrule
\end{tabular}
\end{tabular}
\end{center}

\paragraph{AEMET.} Stabilité : $k=2$, ARI 0.27. nselectboot : $k=4$, ARI 0.39.
Les 4 zones sont mieux séparées avec nselectboot.

\paragraph{Growth.} Les deux ont $k=2$. Stabilité : ARI 0.72 (meilleur $(\alpha,
\omega)$). nselectboot : ARI 0.42. Ici la stabilité fait mieux car elle a choisi
un meilleur couple $(\alpha, \omega)$.

\paragraph{Tecator.} Stabilité : $k=2$, ARI 0.47. nselectboot : $k=3$ (le vrai
$k$), ARI 0.58. Avec la grille fine, nselectboot récupère les 3 classes.

% ============================================================================
\section{Sensibilité au nombre de répétitions $B$}
% ============================================================================

On a testé $B$ = 50, 55, 70, 150. Le triplet optimal change un peu selon $B$,
surtout pour Canadian. Avec $B=150$ et la grille fine, on converge vers $k=2$
partout. Augmenter $B$ ne corrige pas le biais vers $k=2$.

% ============================================================================
\section{Autres méthodes de la littérature}
% ============================================================================

\begin{itemize}
  \item \textbf{Clest} (Dudoit \& Fridlyand, 2002) : split train/test, on
    clustère le train, on prédit sur le test. On compare à des données sans
    structure ($H_0$). Sans cette calibration, notre prototype donne aussi
    $k=2$.
  \item \textbf{Wang (2010)} : validation croisée pour estimer l'instabilité.
    Même idée que Fang-Wang, avec une autre façon de tirer les échantillons.
\end{itemize}

Références en fin de rapport.

% ============================================================================
\section{Pistes pour la suite}
% ============================================================================

\begin{enumerate}
  \item \textbf{Coupler les méthodes} : utiliser la stabilité pour $(\alpha,
    \omega)$ et nselectboot pour $k$.
  \item \textbf{Vote} : faire voter plusieurs critères (stabilité,
    nselectboot, silhouette, gap) et prendre le choix majoritaire.
  \item \textbf{Calibration $H_0$ pour Clest} : implémenter la partie
    « données sans structure » pour voir si Clest fait mieux.
\end{enumerate}

% ============================================================================
\section{Heatmaps : stabilité vs nselectboot}
% ============================================================================

Deux types de heatmaps, car les deux protocoles produisent des grilles
différentes.

\subsection{Stabilité (notre protocole)}
Pour chaque $k$, la stabilité (moyenne ARI) en fonction de $(\alpha, \omega)$.
Rouge = haut = bon. $\times$ = meilleur $(\alpha, \omega)$ pour ce $k$.

\subsection{nselectboot}
Pour chaque $k$, l'instabilité en fonction de $(\alpha, \omega)$. Bleu = bas =
bon (peu instable). $\times$ = meilleur $(\alpha, \omega)$ pour ce $k$.

\subsection{Canadian — Stabilité}
\begin{center}
\includegraphics[width=0.85\textwidth]{figures/stabilite_heatmap_canadian.png}
\end{center}

\subsection{Canadian — nselectboot}
\begin{center}
\includegraphics[width=0.85\textwidth]{figures/nselectboot_heatmap_canadian.png}
\end{center}

\subsection{Autres datasets}
Les figures \texttt{stabilite\_heatmap\_*.png} et
\texttt{nselectboot\_heatmap\_*.png} sont dans \texttt{figures/} pour les 4
datasets. Grille identique : 21×21 (pas 0.05).

% ============================================================================
\section{Paramètres des expériences}
% ============================================================================

\begin{itemize}
  \item \textbf{Stabilité} : grille $\alpha, \omega \in \{0, 0.05, \ldots, 1\}$
    (21×21), $k \in \{2, \ldots, 6\}$, $B=150$, sous-échantillon 80\,\%.
  \item \textbf{nselectboot} : mêmes paramètres (grille 21×21, $B=150$, $k \in
    \{2, \ldots, 6\}$).
\end{itemize}

% ============================================================================
\section{Conclusion}
% ============================================================================

La stabilité est un critère simple et sans labels. Mais elle choisit toujours
$k=2$, ce qui pose problème sur Canadian, AEMET et Tecator. nselectboot, en
testant plusieurs $(\alpha, \omega)$, récupère $k=4$ sur Canadian et AEMET, et
$k=3$ sur Tecator.
Une piste : combiner les deux méthodes ou faire voter plusieurs critères.

\vspace{1em}
\noindent\textbf{En résumé pour les correcteurs :}
\begin{itemize}
  \item Problème : choisir $k$, $\alpha$, $\omega$ sans labels.
  \item Méthode 1 (stabilité) : maximise l'accord quand on enlève 20\,\% des
    données. Résultat : $k=2$ partout.
  \item Méthode 2 (nselectboot) : minimise l'instabilité sur une grille.
    Résultat : $k=4$ sur Canadian et AEMET (correct), $k=2$ sur Growth (correct),
    $k=3$ sur Tecator (correct).
  \item Conclusion : la stabilité a un biais vers les petits $k$. nselectboot
    fait mieux quand on regarde toute la grille.
\end{itemize}

% ============================================================================
\section*{Références}
% ============================================================================

\begin{itemize}
  \item Dudoit, S., Fridlyand, J. (2002). A prediction-based resampling
    method for estimating the number of clusters. \textit{Genome Biology},
    3(7), research0036.
  \item Fang, Y., Wang, J. (2012). Selection of the number of clusters via the
    bootstrap method. \textit{Computational Statistics and Data Analysis},
    56(3), 468--477.
  \item Wang, J. (2010). Consistent selection of the number of clusters via
    crossvalidation. \textit{Biometrika}, 97(4), 893--904.
\end{itemize}

\end{document}
